\section{Operational semantics}
\label{sec:operational-semantics}

\subsection{Values and environments}

Values, environments, outcomes and results are given in \figref{values}. A class value $\Type{C}$ is the
context entry of the class, and an object $\exObj{C}{\bind{\seq{x}}{\seq{v}}}$ records its class by name and fields in the order given by $\fields{\Sigma}(C)$. The value $\Unbound$
marks a variable of the current scope before its declaration is reached, which Python reports as an unbound
local. An environment binds a variable to $\Unbound$ wherever the corresponding context binds the variable to
the entry $\Unbound$, so that a shadowed binding of an enclosing scope is not visible. Throughout this section, $\Sigma$ is the class table produced by module checking of the program
(\secref{well-formedness-modules}) and, like $\mathcal{M}$ and $\mathcal{P}$, is ambient. For a statement $s$, the outcome
$\assigns{\rho}$ represents $s$ running to its end, with accumulated environment $\rho$; the outcome
$\returns{v}$ represents $s$ returning the value $v$; and the outcome $\aborts{\kappa}$ represents the whole run
aborting, with termination kind $\kappa$.

\begin{specdefinition}{Composition of outcomes}
Composing statements sequentially composes their outcomes according to the following monoid, analogous
to sequential composition of static outcomes:
\begin{defalign}
\assigns{\rho} \override \assigns{\rho'} & = \assigns{\rho \override \rho'} \\
\assigns{\rho} \override o & = o \quad (\text{otherwise}) \\
\returns{v} \override o & = \returns{v} \\
\aborts{\kappa} \override o & = \aborts{\kappa}
\end{defalign}

\noindent with $\assigns{\varnothing}$ as left and right unit, and $\returns{v}$ and $\aborts{\kappa}$ as left zeros.
\end{specdefinition}

\begin{figure}
\small
\paperonly{\footnotesize}
\begin{center}
\begin{tabular}[t]{rcll}
\grammarGroup{Values} \\
$v$ & $::=$ & $\ell$ & \speconly{(Literal)} \\
    & $\mid$ & $\exList{v_1, \ldots, v_n}$ & \speconly{(List)} \\
    & $\mid$ & $\exTuple{v_1, \ldots, v_n}$ & \speconly{(Tuple)} \\
    & $\mid$ & $\exDict{\delta}$ & \speconly{(Dict)} \\
    & $\mid$ & $\exLamClosure[\Gamma]{\rho}{\seq{x}}{e}$ & \speconly{(LambdaClosure)} \\
    & $\mid$ & $\exDefClosure[\Gamma]{\rho}{\seq{d}}{i}$ & \speconly{(DefClosure)} \\
    & $\mid$ & $\exObj{C}{\bind{x_1}{v_1} \cons \ldots \cons \bind{x_n}{v_n}}$ & \speconly{(Object)} \\
    & $\mid$ & $\modStub{q}$ & \speconly{(ModuleStub)} \\
    & $\mid$ & $\modLoaded{q}{\rho}$ & \speconly{(LoadedModule)} \\
    & $\mid$ & $\Type{C}$ & \speconly{(Class)} \\
    & $\mid$ & $\TypeAlias[\Gamma]{q}{\seq{\chi}}{i}$ & \speconly{(TypeAlias)} \\
    & $\mid$ & $\Unbound$ & \speconly{(Unbound)} \\
    & $\mid$ & $\primVal{\primFn}$ & \speconly{(Primitive)} \\
    & $\mid$ & $\opaqueVal{q.x}$ & \speconly{(Opaque)}
   \\[3mm]
\grammarGroup{Environments} \\
$\rho$ & $::=$ & $\set{\bind{x_1}{v_1}, \ldots, \bind{x_n}{v_n}}$ & 
   \\[3mm]
\grammarGroup{Dictionary entries} \\
$\delta$ & $::=$ & $\bind{w_1}{v_1} \cons \ldots \cons \bind{w_n}{v_n}$ &
\end{tabular}
\paperonly{\hfill}\speconly{\\[3mm]}
\begin{tabular}[t]{rcll}
\grammarGroup{Outcomes} \\
$o$ & $::=$ & $\assigns{\rho}$ & \speconly{(Assigns)} \\
  & $\mid$ & $\returns{v}$ & \speconly{(Returns)} \\
  & $\mid$ & $\aborts{\kappa}$ & \speconly{(Aborts)}
   \\[3mm]
\grammarGroup{Results} \\
$r$ & $::=$ & $\val{v}$ & \speconly{(Val)} \\
  & $\mid$ & $\aborts{\kappa}$ & \speconly{(Aborts)}
   \\[3mm]
\grammarGroup{Match results} \\
$m$ & $::=$ & $\matchOk{\rho}$ & \speconly{(Match)} \\
  & $\mid$ & $\noMatch$ & \speconly{(NoMatch)} \\
  & $\mid$ & $\kappa$ & \speconly{(Kind)}
   \\[3mm]
\grammarGroup{Termination kinds} \\
$\kappa$ & $::=$ & $\errType \mid \errIndex \mid \errKey {} \mid \errOverflow \mid \errValue$ & \\
    & $\mid$ & $\errZeroDiv \mid \errAssertNoMsg \mid \errAssert{w}$ & \\
    & $\mid$ & $\errAttr \mid \errName \mid {} \errExit{n} \mid \errUndef {} \mid \errExcluded$ &
\end{tabular}

\end{center}
\caption{Values, environments, outcomes and results}
\label{fig:values}
\end{figure}


A \emph{termination
kind} $\kappa$ records how a run ends; all but \replaced{$\errUndef$ and $\errExcluded$}{$\errExit{n}$ and $\errUndef$ } are named after the exceptions raised by the same program
under Python in the same circumstances\added{, except that $\errExit{0}$ also records a run which completes normally}. \Added{A variable read before its assignment aborts with $\errName$ where
Python raises \pyinline{UnboundLocalError}, a subclass of \pyinline{NameError}; distinguishing the two would require
the environment to record whether unbound variables belong to functions or modules.} A run aborts with $\errUndef$ when it reaches a dynamically excluded operation, \added{namely an operation allowed by Python
but not defined by \PurePy,} such as a
comparison of two functions (\secref{no-coercions}). \Added{A run aborts with $\errExcluded$ when it reaches an operation
ruled out by the type system but which Python executes without failing, such as a condition which evaluates to a non-Boolean.
Attribute references and dictionary keys ruled out by the type system also abort with $\errExcluded$, even though Python
sometimes rejects them; to be more precise here would require the semantics to model the attributes Python assigns to a
value, and the keys which Python considers hashable (\secref{eval-expressions}).} A primitive
$\primVal{\primFn}$ is a predefined function whose meaning $\primFn$ is a function from
arguments to results (\specfigref{predefined-values}). Calling $\primVal{\primFn}$ with arguments $\seq{u}$ has result
$\primFn(\seq{u})$. \Added{The meaning of \pyinline{round} aborts with $\errExcluded$ for a Boolean argument, which Python
rounds as the integer 0 or 1, and for a second argument giving the number of digits, and the meaning of
\pyinline{sys.exit} for no argument or an argument other than an integer, for which Python exits with status 0 for
\pyinline{None} and otherwise prints the argument and exits with status 1.}

The rules are given in a monadic metalanguage over results (\figref{values}). If $r$ is
$\val{v}$, then $\mlet{v}{r}{\ldots}$ is the body with $v$ bound to that value; if $r$ is
$\aborts{\kappa}$, it is $\aborts{\kappa}$. The body is a result in the expression rules and an outcome in
the statement rules. A final
premise needs no propagation rule; earlier premises have one each, recording the subexpressions evaluated
before the failing one (\specfigref{eval-propagation}). The sequence and module judgements use
results carrying a sequence of values or an environment in place of a value.
\paperinput{paper/operational-semantics/evaluation-example}

\specinput{spec/operational-semantics/auxiliary}
\paperinput{paper/operational-semantics/operators}

\subsection{Expressions and statements}
\label{sec:eval-expressions}

A value is well-formed at a context entry, and an environment at a context (\specfigref{well-formed-values}).
For any context $\Gamma$, any $\bind{\rho}{\Gamma}$ and any statement $s$ with
$\Gamma \vdash s : R$, the statement evaluation judgement $\rho, s \Rightarrow o$ says that $s$ evaluates
to outcome $o$ in $\rho$ (\specfigrefTwo{eval-block-outcome}{eval-block-fails}). \Added{An \pyinline{if} or
\pyinline{assert} statement aborts with $\errExcluded$ when a condition evaluates to a non-Boolean\speconly{ (\ruleName{eval-if-nonbool}, \ruleName{eval-if-else-nonbool}, \ruleName{eval-assert-nonbool},
\ruleName{eval-assert-msg-nonbool})}, and a failing \pyinline{assert} statement also aborts with $\errExcluded$ when
its message is not a string\speconly{ (\ruleName{eval-assert-msg-nonstr})}.} Top-level statements
are evaluated by a separate judgement, $q, \rho, t \Rightarrow o$ (\specfigref{eval-top-level}), mirroring the static
judgement (\specfigref{typing-top-level}): a sequence composes in the same way as a statement sequence, a mutual type region binds each of its names to $\typeEntry{\Gamma}(q, \seq{\chi}, i)$\speconly{ (\ruleName{eval-types})}, since
class and alias entries are also values, and a plain statement is evaluated by the statement judgement\speconly{
(\ruleName{eval-top-stmt})}. Function definitions delimit
the scope of sequential-composition effects, so as black boxes functions are always pure. The auxiliary function $\result$\speconly{ (\S\ref{aux:result})} reads the outcome of a statement as
the result of a call\speconly{; if the statement completes without returning, the result is $\val{\pyinline{None}}$.}\paperonly{:}
\paperinput{paper/operational-semantics/result}

For any context $\Gamma$, any $\bind{\rho}{\Gamma}$ and any expression $e$ with $\Gamma \vdash e \synth
\tau$, the expression evaluation judgement $\rho, e \Rightarrow r$ (\replaced{\specfigrefTwo{eval-result}{eval-fails}}{\specfigref{eval-result} }) says that $e$ evaluates to
result $r$ in $\rho$, where $r$ is either $\val{v}$ or $\aborts{\kappa}$ (\figref{values}). \Added{A variable
not bound in $\rho$, or bound to $\Unbound$, aborts with $\errName$\speconly{ (\ruleName{eval-var-unknown},
\ruleName{eval-var-unbound})}. A call of a value other than a closure or primitive, a call with the wrong number
of arguments and a constructor call with arguments which do not cover the fields of the class abort with
$\errType$\speconly{ (\ruleName{eval-call-nonfun}, \ruleName{eval-call-lambda-arity}, \ruleName{eval-call-def-arity},
\ruleName{eval-constr-arity})}, and a reference to an attribute absent from an object with $\errAttr$\speconly{
(\ruleName{eval-attr-missing})}, as in Python.} The operators
\pyinline{and} and \pyinline{or} evaluate their second operand only when the first leaves the result open, \Added{aborting with $\errExcluded$ if the first operand is
not a Boolean\speconly{ (\ruleName{eval-and-nonbool}, \ruleName{eval-or-nonbool})},} and
a conditional expression requires its condition to be \replaced{a Boolean}{\pyinline{True} or \pyinline{False} }, like the
\pyinline{if} statement\speconly{ (\ruleName{eval-if})}\Added{, aborting with $\errExcluded$ otherwise\speconly{
(\ruleName{eval-cond-nonbool})}}.

A comprehension evaluates its body once for each binding produced by its qualifiers, and
concatenates the results; a generator draws its bindings from the elements of a value, given by the auxiliary
function $\iterOp$\speconly{ (\S\ref{aux:iter})}: the elements of a
list, tuple or string, and the keys of a dictionary\added{ and the numbers of a range}, as in Python. \Added{A generator over a
value other than a list, tuple, string, dictionary or range aborts with $\errType$\speconly{
(\ruleName{eval-qual-generator-nonseq})}, as in Python, and a guard aborts with $\errExcluded$ if it
evaluates to a non-Boolean\speconly{ (\ruleName{eval-qual-guard-nonbool})}.}
\Added{Dictionary expressions and comprehensions with keys other than strings abort with
$\errExcluded$\speconly{ (\ruleName{eval-dict-nonstr}, \ruleName{eval-dict-comp-nonstr})}, as do subscripting and
membership with such a key; Python raises \pyinline{TypeError} only for keys it cannot hash, such as lists, but \PurePy
has no notion of hashing and cannot distinguish them from other keys.}
\replaced{Subscripting is given by the auxiliary function $\getitem$, and membership by
$\contains{\Sigma}$\speconly{ (\S\ref{aux:getitem}, \S\ref{aux:contains})}; both abort with $\errExcluded$ for a range, as does
subscripting for a Boolean index, which Python treats as the integer 0 or 1.}{Subscripting, given by the auxiliary
function $\getitem$\speconly{ (\S\ref{aux:getitem})}, indexes a list, tuple or
string by an integer, counting from the end when negative, and a dictionary by a string key,
aborting with $\errIndex$, $\errKey$ or $\errType$ when Python raises the exception of that name, and with $\errType$ for a
dictionary subscripted by a value other than a string. } An attribute reference on a
value other than a module or object aborts with \replaced{$\errExcluded$}{$\errAttr$ }\speconly{ (\ruleName{eval-attr-nonobject})}\replaced{. In this version of \PurePy such values have no attributes, whereas in Python a string, for example, has methods; a reference rejected by Python with \pyinline{AttributeError} also aborts with $\errExcluded$, since \PurePy does not keep track of which attributes Python provides}{; this version of \PurePy does not support the attributes given to other values by Python, such as the methods of a string }.

\specinput{spec/fig/well-formed-values}

\specinput{spec/fig/eval-outcome}

\subsection{Pattern matching}
\label{sec:eval-patterns}

A \pyinline{match} statement takes the first case whose pattern matches the value, running its body under
the bindings produced by the pattern, and does nothing if no case matches; the auxiliary function
$\dispatch$\speconly{ (\S\ref{aux:dispatch})} selects the matching case.
\paperinput{paper/operational-semantics/match}%
Whether a pattern matches is determined by the auxiliary function $\matchOp{\Sigma}(\rho, p, v)$\speconly{
(\S\ref{aux:match})}. A \emph{match result} $m$ is
$\matchOk{\rho'}$, a match, where $\rho'$ binds the variables of the pattern to the values that they matched\replaced{;}{, }
$\noMatch$, a failure to match\replaced{;}{, } or \replaced{a termination kind $\kappa$, if matching aborts, such as $\errUndef$}{$\errUndef$, } for a numeric literal pattern if the equality of the literal and the value is undefined. When
\replaced{matching aborts with $\kappa$ before any pattern matches, $\dispatch$ gives $\kappa$, and the run aborts with $\kappa$}{a match result is undefined, the value of $\dispatch$ is also undefined, and the run has undefined behaviour }\speconly{ (\ruleName{eval-match-abort})}.
\speconly{\Added{Matching a list pattern against a tuple, a tuple pattern against a list, or either against a range
aborts with $\errExcluded$ (\S\ref{aux:match}); a constructor pattern against an object aborts with $\errType$ or
$\errExcluded$ when the sub-patterns do not cover the fields of the class (\S\ref{aux:match-fields}).}}
\begin{specdefinition}{Composition of match results}
The match results of sub-patterns compose with
$\override$, matching sequentially, so that the first failure or \replaced{abort determines the result}{undefined result decides }:
\begin{defalign}
\matchOk{\rho} \override \matchOk{\rho'} & {}= \matchOk{(\rho \disjunion \rho')} \\
\matchOk{\rho} \override \noMatch = \noMatch \override m & {}= \noMatch \\
\matchOk{\rho} \override \replaced{\kappa}{\errUndef\ } = \replaced{\kappa}{\errUndef\ } \override m & {}= \replaced{\kappa}{\errUndef\ }
\end{defalign}

\noindent with $\matchOk{\varnothing}$ as left and right unit.
\end{specdefinition}
\paperinput{paper/operational-semantics/pattern-list}%
A constructor pattern matches any object
whose class has the pattern's class among its ancestors, and its sub-patterns bind the fields of the
pattern's class\paperonly{:}\speconly{.}
\paperinput{paper/operational-semantics/pattern-constr}%
Since an object records its class by name, the class test of $\matchOp{\Sigma}$ is a comparison of
fully-qualified names.


\subsection{Modules and programs}

\speconly{Module and program evaluation are given in \figref{eval-modules}; import evaluation in \figref{eval-import-prefix}; importing names from a loaded module in
\figref{eval-imports}.}\paperonly{Module and program evaluation are given in full in \specfigref{eval-modules}.}\speconly{ Context extension (\secref{well-formedness-modules}) applies to environments as well, with the same equations on values in place of context entries, except that a loaded module value $\modLoaded{q}{\rho}$ carries an environment whereas the context entry $\modLoaded{q}{\Gamma}$ carries a context.} \emph{Loading}, the judgement $q \Rightarrow_{\mathsf{mod}} r$, evaluates the import prefix
of $q$ and then its statement, mirroring module checking $\vdash_{\mathsf{mod}} q : \Gamma$
(\specfigref{typing-module}); as with the static rules, the mutual
recursion between loading and evaluation is grounded by program well-formedness, which rules out
import cycles (\secref{modules-programs}).
Loading takes a single module; an import loads the modules containing its target from the root down, each
before the module it contains, mirroring the static rules\speconly{, by the submodules judgement for a plain import and
as premises of the from-import rule}. If loading one of these modules aborts, the import aborts for the same reason and the modules after it are
not loaded.

Since
imports precede statements, every module that a statement can reference is loaded before the statement runs, and
loading is deterministic, so the semantics needs no cache
of loaded modules. Loading runs\speconly{ (\ruleName{eval-module})} a module body in $\mathcal{P}(\pyinline{builtins})$ overridden by
the environment of its imports, with each name declared in the body or bound by a pattern in it starting
unbound, mirroring the \ruleName{module} rule (\secref{well-formedness-modules}).
A program evaluates to a termination kind (\figref{values}): $\errExit{0}$ if the main module
loads, and otherwise the kind it aborted with. $\errExit{n}$ is a run that exited with status $n$\paperonly{:}\speconly{.}
\paperinput{paper/operational-semantics/program}%
An implementation must agree with the semantics on the termination kind of every run with defined behaviour; the form in
which the kind is reported is not prescribed.

\specinput{spec/fig/eval-modules}

\paperinput{paper/operational-semantics/metatheory}
